\section{PF26: the algebraic input and its smaller analytic premise}
\label{sec:pf26-audit}

\subsection{Source and scope of the check}

The source checked here is the supplied 48-page PDF of OpenAI,
\emph{Deterministic Polynomial Factorization over Prime Fields}, dated
4 October 2026.  Printed and PDF page numbers agree. Its SHA-256 is
\texttt{ed88aa5ba802311a21c4f265bcef43515a16019f4b5afdd7f52ecbd36ce0204d}.
The complete proofs in Sections 2--10, pp.~6--44, were read, including all
the inputs to Theorem~1.2, Proposition~8.1 and Proposition~10.1. The
conclusions below are mathematical proof-reading conclusions, not a
formal verification or a claim that an implementation was run. PF26
remains an explicitly named outside preprint in the input register.

\paragraph{MA-PF-001 (the three interfaces are distinct).}
Put
\[
 B=20+(d+1)(\lceil\log_2p\rceil+1).
\]
The algebraic reduction of PF26, Theorem~1.2, p.~3, assumes an auxiliary
prime $\ell_r$ for each prime $r\le d$ only in the branch
$p>B^{200000}$. Its conditions, equation~(1.2), are
\[
 \ell_r\notin\{2,3,p,r\},\qquad
 \ell_r\equiv1\pmod{12r},\qquad
 p^{(\ell_r-1)/r}\not\equiv1\pmod{\ell_r}.
\]
The running time is polynomial in $B$ and the \emph{numerical values}
of these primes. Proposition~10.1, pp.~41--44, gives the explicit
sufficient bound
\[
 O(B^{500000}+B^{100}E^{20}),\qquad E=\max(2,\max_r\ell_r).
 \tag{MA-PF-1}
\]
The more flexible algebraic interface is Proposition~8.1, pp.~35--37:
for a completely split square-free polynomial of odd degree $N$, choose
an odd prime $r\mid N$; supply a field $K/\mathbb F_p$ of degree at
most $r-1$, a primitive $r$th root of unity in $K$, and a generator of
the full $r$-primary subgroup of $K^\times$. These data suffice for a
proper factor in a fixed polynomial number of bit operations. An
auxiliary prime is one way to construct these data; it is not part of
this latter interface.

The analytic premise used for the unconditional statement of PF26 is
\emph{not} the Riemann hypothesis. It is the theorem quoted as
Theorem~11.1, pp.~44--45, from the companion
\emph{Primitive roots for every admissible integer base}, Theorem~1.2:
for every cyclotomic field $K$ containing $\mu_{12}$ and every
finite-order Hecke character $\chi$ of $K$,
\[
 L_K(s,\chi)\ne0\quad\text{if }\Re s>1-10^{-6}.
 \tag{MA-PF-2}
\]
The principal character is included, its pole at $1$ is permitted, and
there is no conductor or imaginary-height restriction. PF26 does not
prove that companion theorem. Its own proof of the implication from
this premise is Lemma~11.2 and Proposition~11.3, pp.~45--46. The
algebraic reduction in Theorem~1.2 uses neither this premise nor zero
location. GRH is a sufficient substitute, as PF26 explicitly observes
on p.~47. In the present paper $A$ is the Dirichlet theorem under the
source convention fixed in the introduction, whereas $M$ is conditional.

\emph{Verification.} The stated distinctions follow directly from the
three source interfaces and the proofs: Sections~2--10 terminate after
constructing the primary table; Section~11 alone bounds the numerical
auxiliary-prime search. In particular, accepting the algebraic PF26
input does not accept its companion's analytic assertion. The proof
reading below found no gap in these algebraic interfaces.\hfill$\square$

\subsection{The completed check of CX-ALG-028}

\paragraph{MA-PF-002 (the finite residual family and an elementary test).}
For the large-characteristic branch $p>B^{200000}$, define
\[
 \mathcal B(p,m,d)=
 \{r\le d:r\text{ prime},\ r\nmid(p-1)m,
              \ p^{r-1}\equiv1\pmod{r^2}\}.
 \tag{MA-PF-3}
\]
Every member is odd and different from $p$. This is exactly the set in
GRH\_NIGHT CX-ALG-028. Under the algebraic PF26 input and $A$, adding
RH for $\zeta_{\mathbb Q(p^{1/r})}$ for $r\in\mathcal B(p,m,d)$
suffices to factor degree-$d$ polynomials over an explicitly supplied
$\mathbb F_{p^m}$ in deterministic polynomial time in
$(d+1)m\log p$. Thus the additional family is empty precisely when
\[
 \forall r\le d\text{ prime},\quad
 r\mid(p-1)m\quad\text{or}\quad p^{r-1}\not\equiv1\pmod{r^2}
 \tag{MA-PF-4}
\]
in that branch. In the small-characteristic branch no additional
analytic input is needed, irrespective of this arithmetic test.
Only primes actually required by the primary-table construction need
be included. This is a sufficient family, not a necessity theorem.

\emph{Proof and independent check of the transfer.}
The large-characteristic inequality gives $p>d$ and $p>3$. For an odd
prime $r\le d$ put $s=\operatorname{ord}_r(p)$ and
$a=v_r(p^s-1)$. Since $s\mid r-1$, writing $r-1=st$ gives
$r\nmid t$. The binomial identity, or the odd-prime lifting-the-exponent
identity, gives
\[
 v_r(p^{r-1}-1)=v_r((p^s)^t-1)=v_r(p^s-1)=a.
\]
Consequently the condition $a>1$ in CX-ALG-028 is precisely the
congruence in (MA-PF-3). The prime $2$ divides $p-1$, so it cannot
belong to the displayed set.

The quadratic primary entry can be obtained from a nonsquare $c$ of
$\mathbb F_p$, by raising $c$ to the odd part of $p-1$. If
$v_2(p-1)=1$, use $c=-1$. If $2\mid m$, compute the quadratic
subfield of the supplied field by solving $z^{p^2}=z$ over
$\mathbb F_p$, and choose $0\ne v$ with $v^p=-v$ there. Then
$c=v^2\in\mathbb F_p$ satisfies $c^{(p-1)/2}=-1$. Otherwise
$A$, with the classical least nonresidue consequence recorded in
GRH\_NIGHT CX-ALG-023, supplies $c$ by a polynomially bounded search.
This checks all quadratic cases, including the supplied-subfield case.

Proceed in increasing order through odd primes $r\le d$. Factoring
$\Phi_r$ has degree $r-1$, so uses only earlier table entries. An
irreducible factor gives
$K_r=\mathbb F_p(\zeta_r)$ of degree
$s=\operatorname{ord}_r(p)\mid r-1$. Put $Q=p^s$ and $R=r^a$.
If $a=1$, the element $\zeta_r$ is already the full primary
generator. If $r\mid p-1$, a small $r$th nonresidue furnished by
$A$ in $\mathbb F_p=K_r$, raised to $(p-1)/R$, is a full primary
generator. Testing whether a candidate is an $r$th power uses the
single modular exponent $(p-1)/r$; no factorization of $p-1$ is needed.

For every remaining case it suffices to construct an explicit
$U_r\simeq\mathbb F_{p^r}$. If $r\mid m$, solve
$z^{p^r}=z$ in the supplied field. Its kernel has dimension $r$ and
its inherited multiplication gives $U_r$. If $r\nmid m$, the
unhandled conditions are exactly $r\in\mathcal B(p,m,d)$.
GRH\_NIGHT CX-ALG-024 supplies $U_r$ in polynomial time from $A$
and RH for the indicated pure field. Its proof was reread in this
check: the factorization of zeta functions, the discriminant bound and
the Bach--Sorenson Theorem~5.1 input agree. One can use the invariant
algebra of PF26 Lemma~3.1, p.~11, to implement its final field
construction. Its auxiliary prime needs congruence modulo $r$ only;
the extra factor $12$ in PF26 equation~(1.2) is unnecessary at this
interface.

Here is the remaining algebraic transfer, with its orders checked.
Because $(r,s)=1$,
\[
 V_r=U_r\otimes_{\mathbb F_p}K_r\simeq\mathbb F_{p^{rs}}
\]
is a field of degree $r$ over $K_r$. Solve the $K_r$-linear equation
$u^Q=\zeta_r u$ for $u\ne0$. Such a solution exists since
$\zeta_r$ has norm one and the map $v\mapsto v^{Q-1}$ has precisely
the norm-one subgroup as its image. As
$v_r(Q^r-1)=a+1$ and $u^{Q-1}$ has order $r$, the $r$-part of
$\operatorname{ord}(u)$ is exactly $r^{a+1}=rR$. Hence
\[
 w=u^{(Q^r-1)/(rR)},\qquad \omega_r=w^r
\]
have exact orders $rR$ and $R$. The subgroup of $V_r^\times$ of
order $R$ is contained in $K_r^\times$, so $\omega_r$ belongs to
$K_r$; its coordinates are recovered by the known inclusion and a
linear solve. All dimensions are at most $r(r-1)$ and all exponent
lengths are polynomial in $r\log p$.

The induction is well founded. PF26 Proposition~2.3, pp.~9--10,
uses the quadratic entry to split a completely split polynomial of
even degree. Proposition~8.1, pp.~35--37, uses precisely the odd
entry just supplied. Its proof receives a field containing a full
primary generator; it does not inspect how that generator was
constructed. PF26 Lemma~9.1 and Proposition~9.2, pp.~38--40,
therefore factor every input below the next table prime with earlier
entries. The table is constructed once; it is not recursively rebuilt.
There are at most $d$ entries and polynomially many factoring calls.
The full reading of the implementation and size proof, reported below,
checks that the exponent is independent of $r$ and of the genus.

For completeness, over the supplied extension field take a square-free
$h$ and form
\[
 C=\mathbb F_{p^m}[X]/(h),\qquad
 C_0=\ker(z\mapsto z^p-z).
\]
If $h$ has $t$ irreducible factors, then $C_0\simeq\mathbb F_p^t$.
The kernel is obtained by $\mathbb F_p$-linear algebra on dimension
$m\deg h$. If $t>1$, a basis supplies a nonscalar $b$. Its minimal
polynomial over $\mathbb F_p$ is a square-free product of distinct
linear factors, of degree between $2$ and $t$. Factor it by the
prime-field procedure. For each root $a$, the gcd
$\gcd(h,b-a)$ selects the irreducible factors on which $b$ has value
$a$, producing a proper split. Degrees strictly decrease under
recursion. Derivative gcds handle multiplicities; if the derivative
vanishes, take coefficient $p$th roots by inverse Frobenius
$c\mapsto c^{p^{m-1}}$ and divide all positive exponents by $p$.
In the branch $p\le B^{200000}$, enumeration of the $p$ scalars in
the same fixed algebra is already polynomial in $B$. These arguments
complete the transfer and the empty-set criterion.\hfill$\square$

\paragraph{Unconditional primary-table promise.}
The stronger conclusion in the last sentence of CX-ALG-028 is also
valid: if every prime $r\le d$ either divides $m$ or has
$v_r(p^{\operatorname{ord}_r(p)}-1)=1$, the displayed construction
uses neither $A$ nor an additional number-field zero-location input.
For $r=2$ use the two elementary quadratic branches above; for odd
$r$ use $\zeta_r$ in the valuation-one branch and the supplied
subfield otherwise. The claim is conditional only on the algebraic
PF26 interface.

\subsection{One Kummer Hecke function for each residual field}

\paragraph{MA-PF-003 (the single-character identity).}
Let $r$ be an odd prime, $p\ne r$ a prime, and put
\[
 F_r=\mathbb Q(p^{1/r}),\quad K_r=\mathbb Q(\zeta_r),\quad
 E_r=K_r(p^{1/r}).
\]
Let $\psi_{p,r}$ be any faithful character of
$\operatorname{Gal}(E_r/K_r)\simeq C_r$, interpreted by class field
theory as its primitive finite-order Hecke character of $K_r$.
Choose the convention in which its unramified value is the Kummer
symbol
\[
 \psi_{p,r}(\mathfrak l)
   =\frac{\operatorname{Frob}_{\mathfrak l}(p^{1/r})}{p^{1/r}}
   =\left(\frac{p}{\mathfrak l}\right)_r
   \quad (\mathfrak l\nmid pr),
 \tag{MA-PF-5}
\]
with a fixed embedding $\mu_r\hookrightarrow\mathbb C$. Then,
including every ramified local factor,
\[
 \boxed{\displaystyle
   \frac{\zeta_{F_r}(s)}{\zeta_{\mathbb Q}(s)}
     =L_{K_r}(s,\psi_{p,r}) .}
 \tag{MA-PF-6}
\]
Each faithful choice gives the same induced $L$-function. In particular,
under $A$, RH for $\zeta_{F_r}$ is equivalent to RH for this single
Hecke $L$-function. The residual family in (MA-PF-3) is therefore an
explicit finite list of finite-order Hecke functions, one for each
$r\in\mathcal B(p,m,d)$.

\emph{Proof.} The Eisenstein argument in CX-ALG-024 gives
\[
 G=\operatorname{Gal}(E_r/\mathbb Q)
       \simeq\mathbb F_r\rtimes\mathbb F_r^\times,
\]
acting on the $r$ roots by affine transformations. Let $P$ denote the
permutation representation on $\mathbb F_r$ and $N$ its translation
subgroup. The character of $P-\boldsymbol1$ is $r-1$ at the
identity, $-1$ at each nonzero translation and $0$ at every element
whose multiplicative part differs from $1$. On the other hand,
$\operatorname{Ind}_N^G\psi_{p,r}$ is zero outside $N$ and its
restriction to $N$ is the sum of all $r-1$ nontrivial characters.
At a nonidentity element of $N$ their sum is $-1$ and at the identity
it is $r-1$. Thus
\[
   P\simeq\boldsymbol1\oplus
            \operatorname{Ind}_N^G\psi_{p,r}.
 \tag{MA-PF-7}
\]
The two sides are actual semisimple representations. Artin induction
formalism gives (MA-PF-6); using inertia-invariant spaces in local
Euler factors proves the equality at the ramified primes as well.
Equivalently, one must use the primitive Hecke character attached to
$E_r/K_r$, not arbitrarily remove every prime above $pr$ from its
$L$-function. Since the character is nontrivial, its Hecke function is
entire. In the open critical strip off the middle line $A$ makes
$\zeta_{\mathbb Q}$ nonzero, so (MA-PF-6) gives the RH equivalence.
This also recovers the already established identity of CX-ALG-024,
\[
 \zeta_{E_r}(s)=\zeta_{K_r}(s)
                     L_{K_r}(s,\psi_{p,r})^{r-1}.
 \tag{MA-PF-8}
\]
No general Artin-holomorphy conjecture is used.\hfill$\square$

\paragraph{Geometric formulation.}
Equation (MA-PF-7) is a decomposition of the Artin motive
$h^0(\operatorname{Spec}F_r)$ into its unit summand and the induction
of the rank-one Kummer motive over $K_r$. On
$\operatorname{Spec}\mathcal O_{K_r}[1/(pr)]$, the latter is a
finite-monodromy rank-one $\ell$-adic sheaf, for $\ell\nmid pr$,
with coefficient field containing $\mu_r$. Its good Frobenius
eigenvalues are roots of unity, hence have weight zero. The local
Euler polynomial is the exterior-power identity
\[
 \det(1-T\operatorname{Frob}\mid V^{I})
  =\sum_j(-T)^j\operatorname{Tr}
           (\operatorname{Frob}\mid\bigwedge^j V^{I}).
\]
This places the identity in the representation-ring and $\lambda$-ring
language; it does not deduce complex zero location from finite-field
weights. For finite-monodromy objects the precise foundational
dictionary is SGA~1, Expos\'e~V, \S7, opening paragraphs, original
pp.~140--141, recomposed p.~115 (arXiv:math/0206203v2, PDF p.~131).
The representation-ring use is the same one recorded in GRH\_NIGHT
CX-ALG-024 and its SGA~6 Expos\'e~V framework reference.

\subsection{A fixed strip is sufficient at the residual primes}

\paragraph{MA-PF-004 (a smaller sufficient analytic condition).}
Fix a number $\delta$ with $0<\delta\le1/2$, independently of
$p,m,d$. Assume the algebraic PF26 input and $A$. In the
large-characteristic branch, it suffices to assume
\[
 L_{K_r}(s,\psi_{p,r})\ne0
       \quad\text{for }\Re s>1-\delta,
       \qquad r\in\mathcal B(p,m,d),
 \tag{MA-PF-9}
\]
to obtain deterministic polynomial-time factorization, with an exponent
allowed to depend on the fixed $\delta$. Thus a proof of full RH for
the residual pure fields is more than this algorithm requires.
The assertion does not permit a different uncontrolled width
$\delta(p,m,d)$ for each input.

\emph{Proof.} For one residual $r$, use $K=\mathbb Q(\zeta_r)$ and
$E=K(p^{1/r})$. Equations (MA-PF-8)--(MA-PF-9) and $A$ place
all nontrivial zeros of both $\zeta_K$ and $\zeta_E$ in
$\Re s\le1-\delta$. The standard smoothed formula used and proved
in PF26 Lemma~11.2, p.~45, works for any fixed such $\delta$.
Indeed, for a fixed nonzero nonnegative
$\varphi\in C_c^\infty((1,2))$, write
$\Phi(s)=\int_0^\infty\varphi(t)t^{s-1}\,dt$ and
\[
 \Psi_D(x)=\sum_{\mathfrak p}\sum_{j\ge1}
       \log N\mathfrak p\,
       \varphi\bigl((N\mathfrak p)^j/x\bigr),
\]
where $\mathfrak p$ runs through the prime ideals of $D$ and
$\Delta_D$ is its absolute discriminant. Mellin inversion and the
smoothed explicit formula give
\[
 \Psi_D(x)=x\Phi(1)-\sum_\rho x^\rho\Phi(\rho)
                       +O_\varphi(\log\Delta_D+[D:\mathbb Q]).
\]
The zeros are counted with multiplicities. Rapid vertical decay of
$\Phi$ and the uniform local zero count give
$\sum_\rho|\Phi(\rho)|\ll_\varphi\log\Delta_D+[D:\mathbb Q]$.
Consequently
\[
 \Psi_D(x)=x\Phi(1)+
 O_\varphi\bigl(x^{1-\delta}
               (\log\Delta_D+[D:\mathbb Q])\bigr)
 \tag{MA-PF-10}
\]
for $D=K,E$. This argument uses the fixed strip, not the middle line.

The discriminant estimate in CX-ALG-024 gives
\[
 [E:\mathbb Q]=r(r-1),\qquad
 \log|D_E|\le r(2r-3)\log r+(r-1)^2\log p.
\]
Thus the degree-plus-log-discriminant quantity for either field is
$O(B^3)$, uniformly. Subtracting $r^{-1}$ times the formula for $E$
from the one for $K$ yields
\[
 \Psi_K(x)-r^{-1}\Psi_E(x)
  =(1-r^{-1})x\Phi(1)+O_\varphi(x^{1-\delta}B^3).
 \tag{MA-PF-11}
\]
Take $x=T_\delta B^{20/\delta}$ with a sufficiently large fixed
$T_\delta$. The error-to-main-term ratio is
$O_\varphi(T_\delta^{-\delta}B^{-17})$.

As in PF26 Proposition~11.3, p.~46, terms of prime-power exponent at
least two or residue degree at least two on the $K$ side contribute
at most
\[
 O_\varphi\bigl([K:\mathbb Q]\sqrt{x}\log^2(2x)\bigr).
\]
The degree-one terms above $p$ or $r$ contribute at most
$O_\varphi([K:\mathbb Q]\log(2x))$. A sufficiently large fixed
$T_\delta$ makes their sum smaller than half the positive main term
in (MA-PF-11). Every remaining degree-one prime ideal of $K$ lies
above a rational prime $\ell\equiv1\pmod r$ with
$\ell\notin\{p,r\}$. If $p$ is an $r$th power modulo $\ell$,
that prime ideal splits completely in $E/K$, and its first-power
contribution cancels exactly in (MA-PF-11). Discarding the remaining
negative terms gives an upper bound. Therefore a noncancelling
degree-one prime must occur, furnishing
\[
 \ell\ll_\delta B^{20/\delta},\qquad
 \ell\equiv1\pmod r,\qquad
 p^{(\ell-1)/r}\not\equiv1\pmod\ell.
 \tag{MA-PF-12}
\]
Search by trial division and modular powering. The invariant algebra
of PF26 Lemma~3.1 constructs $U_r\simeq\mathbb F_{p^r}$ using a
number of operations polynomial in $\ell$ and $\log p$. Apply the
transfer in MA-PF-002. Since $\delta$ is fixed, (MA-PF-12) is a
fixed polynomial bound, and all subsequent costs are polynomial.
No twist by $\mathbb Q(\zeta_{12})$ is needed.\hfill$\square$

The abstract motive and finite-\'etale dictionary is unchanged from
MA-PF-003. The difference is analytic: the prime-detection trace uses
an error $O(x^{1-\delta})$, not necessarily the square-root error.
The finite list (MA-PF-9) remains outside what $A$ alone establishes.
The quoted companion theorem (MA-PF-2) has a different, globally
quantified statement and has not been audited here.

\subsection{The first residual character is already imaginary quadratic}

\paragraph{MA-PF-005 (the cubic instance and the cutoff distinction).}
For $d=3$ in the large-characteristic branch,
\[
 \mathcal B(p,m,3)=\{3\}
 \quad\Longleftrightarrow\quad
       p\equiv-1\pmod9\ \text{ and }\ 3\nmid m.
 \tag{MA-PF-13}
\]
For these instances the only extra Hecke function in MA-PF-003 is
the order-three Kummer character for
\[
 \mathbb Q(\zeta_3,p^{1/3})/\mathbb Q(\zeta_3),\qquad
 \mathbb Q(\zeta_3)=\mathbb Q(\sqrt{-3}).
\]
It is not induced by a character over $\mathbb Q$. Infinitely many
primes meet this congruence and the exact large-characteristic cutoff
for fixed $m$ with $3\nmid m$ and $d=3$.

The concrete prime $p=269$ gives a convenient explicit character:
\[
 269^2-1=72360=3^3\cdot2680,
 \qquad \zeta_{\mathbb Q(269^{1/3})}/\zeta_{\mathbb Q}
       =L_{\mathbb Q(\sqrt{-3})}(s,\psi_{269,3}).
 \tag{MA-PF-14}
\]
It is an example of the arithmetic condition, but is \emph{not} a
case where PF26 needs an analytic assumption: for $d=3$ it has
$B=60$ and $269<60^{200000}$, hence falls in the unconditional
small-characteristic branch.

\emph{Proof.} In the large branch $p>3$. The prime $2$ is excluded.
For $r=3$, exclusion of $3\mid p-1$ leaves $p\equiv-1\pmod3$,
whose order modulo $3$ is $2$. Then
$v_3(p^2-1)>1$ is equivalent to $p\equiv-1\pmod9$.
The other exclusion is exactly $3\nmid m$. The polynomial
$X^3-p$ is Eisenstein at $p$ and has nonsquare discriminant
$-27p^2$, so its splitting field has group $S_3$ and its quadratic
subfield is $\mathbb Q(\sqrt{-3})$. Conjugation by the nontrivial
element of $\operatorname{Gal}(\mathbb Q(\sqrt{-3})/\mathbb Q)$
sends the order-three Kummer character to its inverse. These two
characters differ. A character restricted from $G_{\mathbb Q}$
would be invariant under this conjugation; equivalently a Hecke
character pulled back by the norm would be invariant. Thus this
character does not descend from $\mathbb Q$.

Dirichlet's theorem supplies arbitrarily large primes $p\equiv8$
modulo $9$. For fixed $d=3$, the function
$20+4(\lceil\log_2p\rceil+1)$ grows logarithmically; hence
$p>B^{200000}$ for all sufficiently large $p$. This proves the
infinite-family assertion. Trial division by the primes at most
$\sqrt{269}<17$ verifies primality of $269$, and the displayed
arithmetic verifies the example and its branch.\hfill$\square$

This is the concrete connection to the first non-Dirichlet imaginary
quadratic Hecke case in task~1(d). In the Artin-motive formulation it
is the two-dimensional irreducible summand of the degree-three
permutation motive of $X^3-p$. Holomorphy is already supplied by
Hecke induction; the analytic issue is its zeros. Neither the finite
monodromy nor the weight-zero statement supplies those zeros' location.

\subsection{Record of the full PF26 proof check}

The following checks record what was actually inspected. Every listed
proof was read in its entirety from the supplied PDF's extracted
text. The key interface and analytic-premise pages 35, 44 and 45
were additionally rendered and visually inspected. The source's
classical structure theorems remain cited theorems; their entire
literature was not re-refereed.

\begin{enumerate}
\item \textbf{Lemma~2.1, p.~7.} A nontrivial gcd produced over the
working extension is the product of a subset of $\mathbb F_p$-rational
linear factors, so is over $\mathbb F_p$. Polynomial degrees and
pivots reduce to scalar zero tests. A field inverse is attempted only
after every component is nonzero.
\item \textbf{Lemma~2.2, pp.~8--9.} Distinct-degree decomposition
needs no complete factorization. If $r\mid e$, the Frobenius
eigenspace $v^{Q^{e/r}}=\zeta_rv$ projects nontrivially to every
field component. The moment-curve grid has fewer than $h^2$ forbidden
values; the strict characteristic inequality supplies enough distinct
scalars. The valuation of the order of a nonzero solution is exactly
the full $r$-valuation of $Q^e-1$. Primary logarithms are digit loops,
not loops over the primary subgroup. Promise of an $r$th power makes
each logarithm divisible by $r$.
\item \textbf{Proposition~2.3, pp.~9--10.} Negating a difference
changes its primary logarithm by half the two-primary order. Each
unordered pair therefore receives one orientation. The row counts
cannot all be $(N-1)/2$ when $N$ is even, and $p>N$ preserves the
distinction after reduction modulo $p$.
\item \textbf{Lemma~3.1 and Proposition~3.2, pp.~11--13.} Fixed
functions under the index-$r$ cyclotomic subgroup form an algebra
whose geometric factors are its cosets. The verified residue test
makes Frobenius a transitive $r$-cycle, hence the algebra a field.
The tensor-product field and Frobenius eigenvector have the asserted
primary orders. The only factoring call has degree $r-1$, so the
table/factorization induction has no circular dependency.
\item \textbf{Lemmas~4.1--4.2, pp.~15--16.} Injectivity uses
injectivity of $1-\sigma$ on the infinity lattice, not on its image
in the Jacobian. In the label calculation the function action is
$h(X,Y)\mapsto h(X,\zeta^{-1}Y)$, which gives the sign in
equation~(4.7). Hilbert~90 removes the norm-one function; invariant
divisors decompose into reduced ramification coefficients and a
pullback divisor. Equal labels kill the class because the sum of
the ramification generators is zero.
\item \textbf{Lemma~4.3, Lemma~4.4 and Corollary~4.5,
pp.~17--18.} Frobenius minus one lowers the $\lambda$-filtration
by one; multiplication by $r$ lowers it by $r-1$. The binomial
induction gives the precise extension degree $r^{v_r(N)}$.
The quotient relation's second coordinate shows that fixed pairs
give genuinely fixed Jacobian points. Finite-field norm surjectivity
removes the constant obstruction to descending a divisor class and
its principal function.
\item \textbf{Proposition~4.6, pp.~18--19.} The relation needed
for the next label test is proved only after the previous labels
coincide. The last test would force
$Ne\in\lambda^{1+(r-1)v_r(N)}\Lambda$, contrary to the exact
lattice valuation. No unjustified simultaneous identities at all
levels are used.
\item \textbf{Proposition~5.1 and Lemma~5.2, pp.~20--23.}
Fractional ideals use finite quotients with explicit bounds. Their
inversion, intersections and products reduce to linear conditions.
Characteristic polynomials of multiplication by $X$ compress the
bound and compute pushforward divisors, including multiplicities.
The infinity Vandermonde argument supplies the degree bound for
the complete Riemann--Roch ansatz, and reduction gives absolute
degree at most $2g$.
\item \textbf{Lemma~5.3, pp.~23--24.} At each place a grid
contains a parameter attaining every component's least valuation.
The parameter may depend on the place, which is sufficient to
generate the product ideal; no one parameter must work everywhere.
Determinants compute the component norm without splitting the
coefficient algebra. The grid length and generator lists are
polynomial.
\item \textbf{Lemma~5.4, p.~24.} Binary addition and doubling
of reduced divisor classes stores correcting functions by circuits.
Valuations and nonzero leading coefficients propagate without
expanding the rational function. Since the final divisor has zero
coefficient at a ramification point, the final leading coefficient
is its value there.
\item \textbf{Theorem~6.1 and Lemma~6.2, pp.~25--27.} The norm
promise splits the cyclic algebra only for the existence argument.
The algorithm never uses the unknown splitting or norm solution.
The ramified and unramified local matrix models satisfy the
commutation relation with the stated sign. The local residue-root
promises follow from the norm valuation/residue identities.
Solving $Ve=he$ in a rank-one left ideal gives the norm equation
after $r$ iterations.
\item \textbf{Lemma~6.3 and Lemma~6.4, pp.~28--29.} Completed
orders descend by finite congruence data, glue to
$\operatorname{End}\mathcal V$ and have global sections as claimed.
The splitting proof is over the actual finite field. A maximal
line subbundle has quotient summand degrees no larger than its
degree, making the extension groups vanish. Evaluation gives a
full block triangular algebra. Its trace-pairing kernel is the
strictly triangular radical, whose common column kernel is the
first block; this yields the rank-one fiber idempotent by linear
algebra without diagonalizing a general algebra.
\item \textbf{Section~6.3 and Lemma~6.5, pp.~29--32.} A lift
$b$ of the fiber idempotent has constant characteristic polynomial
$T^{r-1}(T-1)$, so $b^r$ is a rank-one idempotent. Trace-dual basis
bounds control every local coefficient; the finite ansatz contains
all global sections. Precision $3M+3$ suffices after multiplication
by the prescribed powers of local uniformizers. The rational
linear solve has degrees bounded by minors, avoiding expression
tree growth.
\item \textbf{Proposition~7.1, pp.~33--34.} Normalizing the
pushforward function to residue one at a specified ramification
point changes a known norm only by an $r$th power in the constant
field. The norm-solver promise is therefore established. For the
norm-zero divisor $D-\operatorname{div}(a)$, the maximum of cyclic
partial sums solves the divisor equation exactly. Subtracting its
degree times the fixed ramification point gives degree zero and
does not change its $1-\sigma$ image.
\item \textbf{Proposition~8.1, pp.~35--37.} The binomial constant
extension has degree $r^{v_r(N)}\le N$ and supplies the entire
primary group. Every chosen division, rather than merely one
special sequence, is rational over that extension. The simultaneous
sum meets its scalar-grid bound. The reached-test invariant is
exactly Proposition~4.6's invariant. Binary infinity coefficients
are never materialized as large divisors.
\item \textbf{Lemma~9.1 and Proposition~9.2, pp.~38--40.}
The moment-curve separator has fewer than $t^3$ bad parameters
for $t$ components. Characteristic-polynomial square-freeness
detects a separating element, and substitution of a proper factor
selects a proper group of components. Recursive degrees decrease;
derivative-zero and repeated-factor branches preserve multiplicities.
The primary table is built only once.
\item \textbf{Proposition~10.1, pp.~41--44.} The displayed
norm, divisor and circuit dimensions are fixed powers of $B$;
the degree of the curve is not held constant. Moment-curve and
interpolation grids are smaller than $p>B^{200000}$. The
auxiliary algebra is charged by its numerical dimension $\ell-1$,
using elimination without an $\ell$-dependent interpolation grid.
The indexed-array estimate
$O(B^{201000}+B^{41}E^{10})$ squares to a bound covered by
(MA-PF-1). The exponent and the branch cutoffs were checked,
including their dependence on values rather than bit lengths of
auxiliary primes.
\end{enumerate}

\paragraph{Grothendieck's geometric contribution to the algorithm.}
For square-free $h$, $\operatorname{Spec}(\mathbb F_{p^m}[X]/h)$
is a finite \mbox{\'etale} scheme; irreducible factors are its connected
components and primitive idempotents are their indicator functions.
SGA~1, Expos\'e~I, Corollary~7.4, recomposed p.~9 (PDF p.~25),
gives the separable-polynomial criterion; Expos\'e~V, \S7,
recomposed p.~115, gives the finite fundamental-group-set dictionary.
PF26's norm algorithm also makes essential use of the projective-line
splitting theorem, not merely a general effectiveness slogan:
Grothendieck (1957), Th\'eor\`eme~2.1, p.~126, is its published
holomorphic source, and p.~127 discusses the algebraic version over
an algebraically closed field. The form over the original finite
field is proved explicitly in PF26 Lemma~6.3, p.~28. The global
section algebra's parabolic fiber and its rank-one idempotent are
the computational bridge. No archive cote for this new algorithm
has been identified, and none is asserted. The statements do not
turn a general existence theorem in EGA, FGA or SGA into a polynomial
algorithm without the size analysis just checked.
